% ======================================================================
%  CAPÍTULO 1 — INTRODUÇÃO
% ======================================================================
\chapter{Introdução}
\label{cap:introducao}

\section{Contextualização}
\label{sec:contextualizacao}

Os veículos aéreos não tripulados (VANTs) de asas rotativas, em especial os
quadrirrotores, deixaram de ser objeto exclusivo de pesquisa e passaram a
ocupar aplicações cotidianas de inspeção de estruturas, agricultura de
precisão, levantamento topográfico, monitoramento ambiental e logística de
última milha. Na pesquisa em robótica, tornaram-se plataforma padrão, por
permitirem experimentação segura e de baixo custo~\cite{mahony2012}, e os
pilotos automáticos de código aberto, hoje maduros, permitem que outros
pesquisadores validem e ampliem os trabalhos existentes~\cite{ebeid2018}.

Do ponto de vista da Engenharia de Controle e Automação, o quadrirrotor é uma
planta de interesse particular. Ele é \emph{subatuado} --- possui seis graus
de liberdade e apenas quatro entradas independentes, os empuxos dos rotores
---, \emph{não linear}, \emph{acoplado} entre os eixos de atitude e de
translação e \emph{instável em malha aberta}~\cite{bouabdallah2007,beard2008}.
Sem realimentação, a aeronave não voa: todo o comportamento observado em voo é
produto de uma malha de controle executada em tempo real, a algumas centenas
de hertz. Dessa malha, a mais interna e mais rápida é a de \emph{atitude}, que
regula a orientação do corpo no espaço. É a orientação que aponta o vetor de
empuxo, e por isso todas as malhas externas --- de velocidade, de posição, de
navegação --- dependem do desempenho da malha de atitude para funcionar.

Os principais projetos abertos de piloto automático, entre eles o
ArduPilot~\cite{ardupilot2026}, licenciado sob a GPLv3, organizam o controle
em cascata: o comando é refinado sucessivamente da posição para a velocidade,
desta para a aceleração, para a atitude e, por fim, para a velocidade angular,
cujo erro é convertido em esforços de rolagem, arfagem e guinada e repassado à
alocação de empuxo entre os motores. Nas duas malhas internas, o ArduPilot
emprega um controlador proporcional sobre o erro de ângulo seguido de um
controlador proporcional-integral-derivativo (PID) sobre o erro de velocidade
angular, ambos precedidos por um modelador de referência que limita a
aceleração angular comandada~\cite{ardupilotdev,ardupilot2026}.

\section{Problema e motivação}
\label{sec:motivacao}

O PID em cascata é consagrado pelo bom desempenho, pela robustez e pela
facilidade de sintonia empírica~\cite[Cap.~10]{astrom2008}: mais de 95\,\%
dos problemas de controle industrial são resolvidos por controle
PID~\cite[Cap.~1]{astrom2008}. No ArduPilot, ele é
sintonizado por mais de uma década de uso em campo e, em cada aeronave, por um
procedimento automático que excita os eixos em voo e ajusta os ganhos pela
resposta observada. Os ganhos resultantes, contudo, não decorrem de um modelo
dinâmico explícito da aeronave, e a estrutura do controlador permanece a mesma
qualquer que seja a planta. Surge daí a pergunta de pesquisa central deste
trabalho: \emph{um controlador projetado por métodos de controle moderno, a
partir de um modelo dinâmico obtido analiticamente e refinado por
identificação, é capaz de igualar ou superar o desempenho do PID nativo quando
efetivamente embarcado no firmware e submetido às mesmas condições de ensaio?}

Responder a essa pergunta exige mais do que uma simulação em ambiente de
projeto. Um controlador que se comporta bem em um simulador de propósito geral
não é equivalente a um controlador que se comporta bem dentro do laço de tempo
real de um piloto automático, onde a lei de controle convive com período de
amostragem fixo, filtragem dos giroscópios, estimação de estados, saturação
dos atuadores, prioridades da alocação de empuxo e comutação de modos de voo.
O ArduPilot oferece um ponto de entrada para essa avaliação: o
\textit{framework} \codigo{AC\_CustomControl} permite substituir as malhas de
ângulo e de velocidade angular nativas por uma implementação própria em C++,
executada imediatamente antes da biblioteca de motores e comutável em pleno
voo por uma chave do rádio~\cite{ardupilot2026}. A comutação é
metodologicamente valiosa: em um mesmo voo, sob a mesma referência e a mesma
perturbação, é possível medir tanto os controladores propostos quanto a linha
de base.

Há ainda um aspecto que distingue o desempenho aceitável em bancada do
desempenho aceitável em operação: a aeronave real voa sob perturbações. Para a
malha de atitude, o que conta são os momentos: os de assimetrias de massa e
de empuxo, os de rajadas que atingem os rotores de forma desigual e os da
perda parcial de um motor (Seção~\ref{sec:perturbacoes}). Avaliar um
controlador apenas pela resposta ao degrau em condição nominal diz, portanto,
pouco sobre sua utilidade prática. No simulador usado neste trabalho, porém,
o vento quase não gera momento de atitude: ele age como força no centro de
massa e desloca a aeronave (Seção~\ref{sec:res-vento}). Por isso, as
perturbações de atitude são aplicadas diretamente como torque, e tratadas ao
mesmo tempo como \emph{cenário de ensaio}, por meio de perfis reprodutíveis
de torque aplicados igualmente a todos os controladores, e como \emph{objeto
de projeto}, por meio de mecanismos de estimação e compensação de perturbação
incorporados à lei de controle. A capacidade de rejeição de perturbação passa,
assim, de efeito colateral a critério explícito de comparação.

A validação é feita em ambiente \textit{software-in-the-loop} (SITL), no qual
o firmware efetivamente compilado é executado contra um modelo de dinâmica de
voo, sem \textit{hardware}~\cite{ardupilotdev}. Os ensaios são plenamente
repetíveis, as perturbações são programáveis e não há risco material ou à
integridade física. Sob a ótica da Engenharia de Controle e Automação, o
trabalho percorre o ciclo completo de projeto de um sistema de controle ---
modelar a planta, identificá-la, projetar a lei de controle, discretizá-la,
embarcá-la e avaliá-la experimentalmente --- com custo material praticamente
nulo.

\section{Objetivos}
\label{sec:objetivos}

\subsection{Objetivo geral}

Projetar, implementar e validar em simulação um controlador de atitude
baseado em modelo para um quadrirrotor, em três variantes (LQI, LQR com
observador de estado estendido e LQG), embarcado no firmware ArduPilot por
meio do
\textit{framework} \codigo{AC\_CustomControl}, comparando quantitativamente
seu desempenho ao do controlador PID nativo em ambiente SITL, com uso
exclusivo de ferramentas de código aberto.

\subsection{Objetivos específicos}

\begin{alineas}
    \item Obter o modelo dinâmico não linear do quadrirrotor a partir das
          equações de Newton-Euler, incluindo a cinemática de atitude, o modelo
          dos rotores e a matriz de alocação, e levantar os parâmetros físicos
          correspondentes ao modelo simulado;
    \item linearizar o modelo em torno do voo pairado, obter as
          representações em espaço de estados contínua e discreta, verificar a
          controlabilidade e a observabilidade e validar a linearização contra
          o modelo não linear;
    \item caracterizar experimentalmente o controlador PID nativo do
          ArduPilot e estabelecê-lo como linha de base quantitativa;
    \item identificar a dinâmica de atitude da aeronave simulada no domínio da
          frequência e confrontá-la com o modelo analítico;
    \item projetar o controlador de atitude pelo regulador linear quadrático
          (LQR) com ação integral (LQI), pelo LQR com observador de estado
          estendido e pelo LQG, verificar as margens de estabilidade e
          projetá-lo em tempo discreto, na taxa de execução do laço;
    \item modelar e implementar perturbações externas reprodutíveis ---
          perturbações de torque em pulso, em degrau sustentado, em varredura
          de frequência e estocásticas, e uma rajada de vento, para verificar
          o seu efeito sobre a atitude no simulador ---, compondo cenários de
          ensaio parametrizáveis, e cenários de robustez com erro de modelo,
          falha de motor, ruído de giroscópio, saturação e grandes ângulos;
    \item estimar e compensar a perturbação, por um observador de estado
          estendido e por um filtro de Kalman, e comparar essa estratégia à
          da ação integral;
    \item implementar os controladores em C++ no ArduPilot, com ganhos expostos
          como parâmetros ajustáveis em tempo de execução, registro dos sinais
          internos e reinicialização adequada dos estados na comutação entre
          controladores;
    \item comparar os controladores em SITL por métricas objetivas de
          desempenho transitório, erro de rastreamento, rejeição de
          perturbação e esforço de controle;
    \item documentar o procedimento de forma reprodutível, com código, cenários
          e análises versionados.
\end{alineas}

\section{Justificativa}
\label{sec:justificativa}

A relevância do tema é reforçada por trabalhos consolidados na área.
\citeonline{bouabdallah2007} desenvolveu a modelagem, o projeto e o
controle de um quadrirrotor, comparando estratégias lineares e não lineares em
uma plataforma experimental construída para esse fim.
\citeonline{sabatino2015}, por sua vez, derivou o modelo não linear de seis
graus de liberdade e projetou controladores por realimentação de estados e
por técnicas não lineares, avaliando-os exclusivamente em simulação numérica.

Este trabalho fica entre essas duas abordagens. Em vez de
construir uma aeronave própria, o que desloca o esforço para a montagem
mecânica e eletrônica, ou de encerrar a avaliação no ambiente de projeto, o
que deixa em aberto a questão da implementação, a proposta consiste em
\emph{embarcar o controlador projetado em um piloto automático de uso real e
de código aberto} e confrontá-lo diretamente com o controlador de produção sob
condições de ensaio idênticas. O controlador passa a enfrentar a mesma cadeia
de estimação, o mesmo modelador de referência, a mesma alocação de empuxo e o
mesmo período de amostragem que o PID nativo enfrenta, e a comparação deixa de
depender de simplificações feitas apenas de um dos lados.

Um resultado adicional esperado é de natureza metodológica: um roteiro
reprodutível, em português, que vai da equação de Newton-Euler até o código
embarcado e até a bateria de ensaios automatizada, construído exclusivamente
com software livre e passível de continuidade por outros alunos.

\section{Trabalhos relacionados}
\label{sec:relacionados}

Três grupos de trabalhos cercam este. O primeiro compara o PID com o controle
ótimo no próprio quadrirrotor. \citeonline{bouabdallah2004} aplicaram ao OS4
um PID, projetado sobre uma dinâmica simplificada, e um LQ, projetado sobre um
modelo mais completo. Em voo, o PID estabilizou a atitude sob perturbações
pequenas, e o LQ teve resultado apenas mediano, que os autores atribuíram às
imperfeições do modelo. A vantagem de um projeto baseado em modelo depende,
portanto, de o modelo representar o que o controlador de fato enfrenta. É por
isso que este trabalho identifica o modelo no próprio SITL
(Capítulo~\ref{cap:modelagem}).

O segundo grupo embarca controladores em pilotos automáticos de código aberto.
O controlador de ângulo do PX4 implementa a lei em quatérnios de
\citeonline{brescianini2013}, segundo a documentação do próprio PX4
\cite{px4controladores}. \citeonline{smeur2016} levaram ao voo de
microveículos a inversão dinâmica incremental adaptativa. Três trabalhos
comparam o controlador proposto com o de produção dentro do mesmo
\textit{firmware}:
\begin{itemize}
\item \citeonline{nguyen2015} compararam, no ArduPilot, o PID original, um
PID novo e um \textit{backstepping} integral;
\item \citeonline{bressan2020} trocaram o controlador de taxa do PX4 por um
controlador adaptativo com observador, que pode ser ligado e desligado em
voo. Eles o compararam com a cascata original sob um torque de perturbação no
pairado, medindo o tempo até a atitude nivelar, o pico do desvio e o esforço
de controle, as mesmas famílias de métrica usadas aqui;
\item \citeonline{saldiran2024} implementaram a inversão dinâmica incremental
como biblioteca à parte no ArduPilot, com chaveamento para o controlador
original. Eles a verificaram no SITL e a compararam em voo com o PID ajustado
pelo AutoTune, a mesma linha de base deste trabalho.
\end{itemize}

O terceiro grupo estima a perturbação e a cancela. O observador de estado
estendido do LQR+ESO vem do controle por rejeição ativa de perturbações de
\citeonline{han2009}, e \citeonline{chen2016} revisam o controle baseado em
observador de perturbação. \citeonline{sanz2016} aplicaram à atitude de um
quadrirrotor um estimador de incerteza e perturbação em que um só parâmetro
troca rejeição por amplificação do ruído de medida no sinal de controle; no
voo, escolheram um valor com esforço próximo ao de um PID de dois graus de
liberdade. \citeonline{lyu2018} projetaram o observador por síntese
$H_\infty$ e notaram que a banda do controlador de base da aeronave deles
era limitada pelo ruído e pelos modos flexíveis.

Esses trabalhos medem o ganho na rejeição. O esforço de controle aparece como
uma métrica a mais \cite{bressan2020}, e o compromisso com o ruído de medida
aparece como verificação depois do projeto \cite{sanz2016}. Este trabalho faz
desse compromisso uma restrição do projeto: um controlador proposto só é
comparado ao nativo se não passar aos motores mais ruído do giroscópio do
que ele (Seção~\ref{sec:projeto-lqg}). Com essa regra, parte da vantagem que
um projeto de banda maior mostraria se revela efeito do ruído, e não da
estrutura do controlador. Além disso, como o LQG age sobre o caminho de medida, o
desempenho é conferido também na atitude verdadeira do simulador, e não só
na estimada pelo AHRS (Apêndice~\ref{ap:bases}).

\section{Delimitação do escopo}
\label{sec:delimitacao}

O trabalho restringe-se à malha de atitude --- rolagem, arfagem e guinada ---
de um quadrirrotor em configuração X, no modelo de multirrotor nativo do SITL
do ArduPilot, com os parâmetros físicos descritos na
Seção~\ref{sec:sitl}. As malhas de altitude, de velocidade e de posição
permanecem as nativas. Toda a validação ocorre em simulação; o firmware
utilizado é a versão de desenvolvimento 4.8 do ArduCopter, fixada em um
\textit{commit} específico do repositório oficial~\cite{ardupilot2026}, de modo
que os resultados possam ser reproduzidos.

O SITL simplifica a aeronave em pontos que favorecem os controladores
baseados em modelo: não tem dinâmica de motor, não modela o termo giroscópico
$\vetor{\omega}\times\matriz{J}\,\vetor{\omega}$, trata o vento quase só como
perturbação de posição, aplica a carga suspensa e o cabo sem momento e, nos
ensaios-base, entrega um giroscópio praticamente sem ruído
(Seções~\ref{sec:sitl} e~\ref{sec:motor-sitl}). O último capítulo discute o
que isso implica para uma aeronave real.

Quanto às perturbações, o escopo abrange perturbações aditivas de torque nos
três eixos --- em pulso, em degrau, em varredura e estocásticas ---, uma
rajada de vento uniforme no volume ocupado pela aeronave e variações da
própria aeronave e dos sensores: erro de modelo, falha de motor e ruído de
giroscópio. O injetor de perturbações implementa também o vento médio e a
turbulência de Dryden (Apêndice~\ref{ap:vento}), que não são usados nos
ensaios. Ficam
fora do escopo os campos de vento espacialmente distribuídos, as esteiras de
obstáculos e os efeitos aerodinâmicos de proximidade ao solo ou a paredes, que
exigiriam simulação fluidodinâmica. Ficam igualmente fora do escopo o voo em
\textit{hardware} real, o controle de trajetória por métodos preditivos, a
operação cooperativa de múltiplas aeronaves e o uso de visão computacional,
registrados como possibilidades de trabalhos futuros.

\section{Organização do trabalho}
\label{sec:organizacao}

Este trabalho está organizado em sete capítulos. O
Capítulo~\ref{cap:fundamentacao} apresenta a fundamentação teórica: a
representação da atitude, o modelo dinâmico do quadrirrotor e sua
linearização, a arquitetura de estimação e de controle de atitude do
ArduPilot, as técnicas de controle empregadas --- PID, LQR, observador de
estado estendido e LQG ---, a integral de Bode, os modelos de perturbação, a
identificação de sistemas no domínio da frequência e o simulador SITL. O
Capítulo~\ref{cap:metodos} descreve os materiais e métodos: o ambiente de
simulação, o injetor de perturbações, as dezessete configurações de
controlador, os 55 cenários de ensaio, as métricas de comparação e a
automação. O Capítulo~\ref{cap:modelagem} trata da modelagem e da
identificação da aeronave simulada e valida o modelo linear contra o SITL. O
Capítulo~\ref{cap:projeto} projeta o LQI, o LQR com observador de estado
estendido e o LQG, descreve a sua implementação no firmware e confere o
código embarcado contra o modelo. O Capítulo~\ref{cap:resultados} verifica o
caminho de integração, caracteriza o PID nativo e compara com ele os
controladores propostos em todos os grupos de cenários. O
Capítulo~\ref{cap:conclusoes} responde à pergunta de pesquisa, discute os
limites encontrados e as limitações do simulador, situa outras técnicas de
controle e propõe trabalhos futuros. O Apêndice~\ref{ap:vento} reúne os
modelos de vento e a sua implementação no injetor de perturbações; o
Apêndice~\ref{ap:validacao-proposto} lista, ensaio a ensaio, o confronto entre o modelo
e o SITL dos controladores propostos; o Apêndice~\ref{ap:reproduzir} descreve
como reproduzir os resultados; e o Apêndice~\ref{ap:bases} compara a
atitude do AHRS com a atitude verdadeira do simulador.
